\chapter{Dicas e Boas Práticas} \label{chp:ap_dicas}

Esse apêndice tem como objetivo explicar como construir um texto científico, apresentando dicas na construção do texto e demonstrando exemplos de códigos em LaTeX para as principais necessidades na elaboração do texto.

% ---
\section{Inclusão de outros arquivos}
% ---

É uma boa prática dividir o seu documento em diversos arquivos, 
e não
apenas escrever tudo em um único. 
Esse recurso foi utilizado neste documento. 
Para incluir diferentes arquivos em um arquivo principal, 
de modo que cada arquivo incluído fique em uma página diferente, utilize o comando:

\begin{verbatim}
   \include{documento-a-ser-incluido}      % sem a extensão .tex
\end{verbatim}

Para incluir documentos sem quebra de páginas, utilize:

\begin{verbatim}
   \input{documento-a-ser-incluido}      % sem a extensão .tex
\end{verbatim}

% ---
\section{Citações}
% ---

Ao redigir textos científicos, é importante utilizar corretamente as diferentes formas de citação para garantir clareza, fluidez e adequação às normas. Uma citação comum ocorre ao final da frase, quando se apresenta uma informação e se associa a referência entre parênteses, como em: A adoção de sistemas híbridos melhora a acurácia das recomendações \cite{braida2015transforming}.

Quando todo o conteúdo de um parágrafo deriva de uma mesma fonte, é possível posicionar a citação apenas ao final do parágrafo: Os sistemas de recomendação evoluíram para incorporar múltiplas fontes de dados, como perfil do usuário, histórico de interações e contexto de uso. Essa combinação visa aumentar a relevância dos resultados gerados, superando limitações dos métodos tradicionais. \cite{braida2015transforming}

Também é comum usar o autor como parte integrante da frase, especialmente quando se deseja destacar a contribuição de um pesquisador. Nesse caso, utiliza-se o comando \verb|\citeonline|. É fundamental prestar atenção à concordância verbal: se o autor for único, o verbo deve estar no singular, como em ``\citeonline{abntex2classe} propõe uma abordagem baseada em fatores latentes''; se houver múltiplos autores, o verbo vai para o plural, como em ``\citeonline{braida2015transforming} propuseram um modelo robusto para classificação supervisionada'' .

Em situações que exigem a reprodução literal das palavras do autor, utiliza-se uma citação direta, que pode ser feita com aspas ou com o ambiente \verb|quote|. Por exemplo:

\begin{quote}
``A combinação de técnicas colaborativas e baseadas em conteúdo tem se mostrado eficaz para sistemas de recomendação personalizados'' \cite{braida2015transforming}.
\end{quote}

Quando o autor é apresentado como sujeito da primeira frase do parágrafo, e todas as ideias subsequentes derivam logicamente da mesma fonte, não há necessidade de repetir a citação a cada sentença. Por exemplo:

\citeonline{braida2015transforming} analisam os principais desafios enfrentados pelos sistemas de recomendação híbridos. Os autores destacam a complexidade de combinar diferentes abordagens sem aumentar significativamente o custo computacional. Além disso, apontam que a personalização pode ser prejudicada se os dados contextuais não forem bem integrados. Por fim, propõem um modelo que equilibra desempenho e escalabilidade, aplicável a diversos domínios.

Essas diferentes formas de citação devem ser usadas estrategicamente ao longo do texto, promovendo coesão e evitando repetições desnecessárias. Além disso, respeitar a concordância entre autor e verbo demonstra atenção aos detalhes e contribui para a qualidade da escrita acadêmica.

Por fim, é importante lembrar que há casos que requerem atenção especial, como quando o trecho citado foi originalmente escrito em outro idioma e o autor da monografia opta por traduzi-lo. Nessas situações, deve-se indicar que a tradução é livre ou do próprio autor (por exemplo, tradução nossa ou tradução do autor), conforme as regras da ABNT. Diante de situações específicas como essa, é sempre recomendável consultar diretamente a norma vigente ou o manual de normalização da instituição para garantir a conformidade formal do trabalho.

% ---
\section{Imagens e Tabelas}\label{sec:LABEL_CHP_2_SEC_A}
% ---

Toda tabela inserida em um trabalho deve obrigatoriamente ser referenciada e explicada no corpo do texto. Não é aceitável incluir uma tabela isoladamente, sem que o leitor seja orientado sobre seu conteúdo, sua relevância ou os dados que ela apresenta.

A referência à tabela deve ocorrer de forma clara e integrada ao parágrafo. Por exemplo:

A Tabela~\ref{tab:MATRIX_NOTAS_EXEMPLO} e a Imagem~\ref{fig:LABEL_FIG_1} resumem as médias obtidas pelos grupos ao longo da disciplina, permitindo uma comparação direta entre os diferentes métodos aplicados.

Note que, no exemplo acima, foi utilizado o símbolo \verb|~| entre a palavra \verb|Tabela| e o comando \verb|\ref|. Isso garante que o LaTeX não quebre a linha entre esses dois elementos, mantendo a referência legível e correta. O mesmo princípio vale para figuras, seções e outros elementos numerados: recomenda-se sempre usar o ~ para manter a unidade da referência.

Além disso, é recomendável evitar expressões genéricas como ``a tabela acima'' ou ``a tabela abaixo'', pois a posição pode variar com alterações no layout. O uso do comando \verb|\ref{}| é a forma correta e robusta de referenciar tabelas, figuras e equações no LaTeX.

No LaTeX, ao inserir uma tabela ou figura, é necessário informar ao compilador sugestões de posicionamento por meio de letras dentro de colchetes, como em  \verb|\begin{table}[!htb]| . Nesse caso, o  \verb|h| indica que o ideal é posicionar o elemento no local exato onde foi declarado (\textit{here}),  \verb|t| permite o topo da página (\textit{top}) e  \verb|b| permite o rodapé (\textit{bottom}). O símbolo de exclamação  \verb|!| serve para flexibilizar as restrições internas do LaTeX, aumentando as chances de o elemento ser posicionado conforme desejado. O uso de  \verb|!htb| costuma produzir bons resultados, equilibrando controle e harmonia visual. No entanto, há situações em que é necessário que a tabela ou figura apareça exatamente onde foi inserida no código, sem qualquer movimentação. Nesses casos, pode-se utilizar o modificador  \verb|[H]| (com  \verb|H| maiúsculo), fornecido pelo pacote  \verb|float|, que força o posicionamento exato no ponto de inserção. Essa abordagem deve ser usada com cautela, pois pode afetar a estética e a organização do documento.


\begin{table}[!htb]
  \centering
  \caption{Fragmento de uma matriz de notas de um sistema de recomendação de filmes.}
  \label{tab:MATRIX_NOTAS_EXEMPLO}
  \rowcolors{1}{}{lightlightgray}
  \begin{tabular}{l c c c}
  \toprule
    & Titanic & Poderoso Chefão & Matrix \\
    \midrule
        Filipe Braida & 4 & $\emptyset$ & 3  \\
        Leandro Alvim & 4 & 5 & 5  \\
        Bruno Dembogurski & 4 & 5 & 5  \\
        Fellipe Duarte & $\emptyset$ & 5 & $\emptyset$  \\ 
    \bottomrule
  \end{tabular}
\end{table}

\begin{figure}
  \centering
  \includegraphics[width=0.6\textwidth]{imagens/chick.png}
  \caption{Chick}
  \label{fig:LABEL_FIG_1}
  \legend{Fonte: o autor.}
\end{figure}

\begin{algorithm}[H]
\floatname{algorithm}{Algoritmo}

\textbf{Entrada} conjunto de notas $R$, limiar $\tau$, modelo de previsão $\varphi$.

\textbf{Saída} conjunto de notas sem ruídos $R^{*}$.

$R^{*} \gets \{\}$

\For{$(u, i, r) \in R$}{
$\tilde{r} \gets \varphi(u,i)$\;

\If{$|\tilde{r} - r| < \tau$}{
$R^{*} \gets R^{*} \cup \{(u,i,r)\}$
}
}
\caption{Filtragem das avaliações com ruído proposto por \citeonline{braida2015transforming}.}
\label{alg:mahony}
\end{algorithm}

\section{Equações}

Equações em textos científicos devem ser tratadas com o mesmo cuidado dedicado a figuras e tabelas. Isso significa que toda equação inserida no documento precisa ser referenciada e explicada no corpo do texto, jamais aparecendo isoladamente ou sem contexto. Assim como escrevemos ``A Figura~\ref{fig:LABEL_FIG_1} mostra...'' ou ``A Tabela~\ref{tab:MATRIX_NOTAS_EXEMPLO} apresenta...'’, também é necessário dizer, por exemplo, ``A equação~\ref{eq:LABEL_EQ_1} descreve...''. Como mostra a equação~\ref{eq:LABEL_EQ_1}, a identidade fundamental do quadrado da soma pode ser expressa da seguinte forma:

\begin{equation}
  (x + y)^2 = x^2 + 2xy + y^2
  \label{eq:LABEL_EQ_1}
\end{equation}

Uma equação pode ser incorporada ao texto de diferentes maneiras. Quando for curta e direta, pode ser inserida no próprio parágrafo, delimitada por cifrões simples (\verb|$|), como em: a fórmula da área do círculo é $A = \pi r^2$. Essa modalidade é conhecida como equação em linha.

Por outro lado, quando se deseja evidenciar a equação em uma linha destacada, utiliza-se o ambiente \verb|equation|, como no exemplo a seguir:
\begin{equation}
  (x + y)^2 = x^2 + 2xy + y^2 \: ,
  \label{eq:LABEL_EQ_2}
\end{equation}
onde $x$ e $y$ representam variáveis reais quaisquer.

Embora apareça destacada visualmente, essa equação ainda faz parte do parágrafo e, portanto, deve seguir a mesma pontuação e estrutura do texto. Isso significa que a frase deve continuar normalmente após a equação, com vírgula, ponto final ou qualquer outro sinal de pontuação adequado.

É importante reforçar que nenhuma equação deve ser apresentada sem uma explicação clara. O leitor precisa compreender o seu propósito, a interpretação dos símbolos utilizados e a sua relação com a argumentação desenvolvida no texto. Equações bem contextualizadas contribuem para a clareza, a precisão e a força do raciocínio científico.



\section{Listings}\label{sec:LABEL_CHP_2_SEC_D}
Reference: \url{http://en.wikibooks.org/wiki/LaTeX/Source_Code_Listings}

\codec{C}{alg:LABEL_CODE_1}{codigos/codigo-c.txt}

\codejava{Java}{alg:LABEL_CODE_2}{codigos/codigo-java.txt}

\section{Referências Cruzadas}\label{sec:LABEL_CHP_2_SEC_E}

Em documentos acadêmicos escritos em \LaTeX, é altamente recomendável utilizar referências cruzadas automáticas para manter a consistência na numeração e facilitar a leitura. Para isso, utiliza-se o sistema de marcação com rótulos (\verb|\label|) e comandos específicos para cada tipo de elemento.

Embora seja possível referenciar manualmente com comandos como \verb|Tabela~\ref|, o uso de comandos específicos torna o processo mais organizado e facilita alterações em todo o documento, especialmente quando se deseja padronizar os prefixos como ``Capítulo'', ``Seção'' ou ``Equação'' automaticamente.

Para capítulos, utiliza-se o comando \verb|\refchapter|, que gera automaticamente o número correspondente do capítulo referenciado. Da mesma forma, para seções, o comando \verb|\refsection| insere o número da seção associada ao rótulo definido.

Tabelas devem ser referenciadas com o comando \verb|\reftable|, e figuras com \verb|\reffigure|. Esses comandos inserem os números correspondentes às legendas das respectivas tabelas e figuras, respeitando a numeração automática do documento.

No caso das equações, utiliza-se \verb|\refequation|, garantindo que a equação seja mencionada corretamente no texto. Para algoritmos, especialmente ao utilizar o pacote \verb|algorithm2e|, a referência deve ser feita com \verb|\reflisting|.

Por fim, quando for necessário referenciar apêndices, pode-se usar \verb|\refappendix|, mantendo a padronização também nas seções posteriores do trabalho.

Esses comandos não apenas facilitam a manutenção do documento como também evitam erros de numeração, garantindo que todas as referências se atualizem automaticamente a cada compilação.

\section{Definições, Teoremas}

\begin{definition}\label{def:def1}
  Aqui é uma nova definição.
\end{definition}

\begin{definition}[Título] \label{def:def2}
  Aqui é uma outra definição.
\end{definition}

\begin{theorem}\label{the:the1}
  Aqui é um teorema.
\end{theorem}

Seguindo a Definição~\ref{def:def1} e Teorema~\ref{the:the1}.



\section{Acrônimos, Siglas}
  Um \ac{SR}... Portanto, o \ac{SR}...


\section{Símbolos}
  A \ac{rv} é dada por... Assim, \ac{rv}...
